\documentclass[11pt]{article}

\usepackage[margin=1in]{geometry}
\usepackage{amsmath,amssymb,amsthm,mathtools}
\usepackage{graphicx}
\usepackage{xcolor}
\usepackage{enumitem}
\usepackage{microtype}
\usepackage[hidelinks]{hyperref}

\newtheorem{theorem}{Theorem}
\newtheorem{lemma}[theorem]{Lemma}
\newtheorem{proposition}[theorem]{Proposition}
\theoremstyle{remark}
\newtheorem{remark}[theorem]{Remark}

\newcommand{\R}{\mathbb R}
\newcommand{\Waux}{\mathcal W}
\newcommand{\Lip}{\operatorname{Lip}}
\newcommand{\mean}{\mathop{\rlap{\raisebox{0.08ex}{\scriptsize--}}\!\int}}

\title{A block-projection counterexample for\\complex interpolation of weighted Sobolev spaces}
\author{Autonomous proof attempt for run \texttt{fa\_banach\_001}}
\date{13 August 2026}

\begin{document}
\maketitle

\begin{center}
\fcolorbox{black}{gray!8}{\parbox{0.89\textwidth}{
\textbf{Status.} Candidate full affirmative answer to the first question in
Section~8.3 of Cwikel--Einav, together with a stronger counterexample to the
unrestricted Stein--Weiss-type identity.  The proof is analytic and is
submitted for human review.
}}
\end{center}

\begin{abstract}
We construct positive compact-bounded weights $w_0,w_1$ on $\R$ for which
$r=w_0/w_1$ is locally Lipschitz and, at $p=2$ and $\theta=1/2$,
\[
 \Waux^2(\R,1/2,r)
 \subsetneq [W^{1,2}(\R,w_0),W^{1,2}(\R,w_1)]_{1/2}
 \subsetneq W^{1,2}(\R,\sqrt{w_0w_1}).
\]
The first strict inclusion answers the exact open question.  The second gives
a counterexample to the unrestricted weighted-Sobolev analogue of the
Stein--Weiss formula.  The mechanism is a common complemented block subspace:
its endpoint norm is equivalent to unweighted $\ell^2$, while its
geometric-mean norm is exponentially smaller.
\end{abstract}

\section{Source question}

The source is M.~Cwikel and A.~Einav, \emph{Interpolation of weighted Sobolev
spaces}, arXiv:1807.10790, later published in the \emph{Journal of Functional
Analysis} 277 (2019), 2381--2441 \cite{CE}.  On PDF page~47, at the start of
Section~8.3, the authors ask:
\begin{quote}
Can one find an open set $U$ and weights $w_0,w_1$ on it for which
\[
 \Waux^p(U,\theta,r_w)
 \underset{\neq}{\subset}
 [W^{1,p}(U,w_0),W^{1,p}(U,w_1)]_\theta ?
\]
\end{quote}

\begin{figure}[ht]
  \centering
  \includegraphics[width=\textwidth]{figures/open_problem_crop.png}
  \caption{The question as rendered from PDF page~47 of arXiv:1807.10790.}
\end{figure}

For positive weights satisfying the source paper's compact boundedness
condition, write
\[
 X_j=W^{1,2}(\R,w_j),\qquad
 r=\frac{w_0}{w_1},\qquad w_{1/2}=\sqrt{w_0w_1}.
\]
The auxiliary space in the question is
\[
 \Waux^2(\R,1/2,r)
 =\left\{u\in W^{1,2}(\R,w_{1/2}):
       u\,| (\log r)'|\in L^2(\R,w_{1/2})\right\}.
\]
Theorem~1.21 of \cite{CE} says that on $U=\R$, if $w_0,w_1$ are
compact-bounded and $r$ is locally Lipschitz, then
\begin{equation}\label{eq:source-inclusion}
 \Waux^2(\R,1/2,r)\subset [X_0,X_1]_{1/2}
 \subset W^{1,2}(\R,w_{1/2}).
\end{equation}

\section{Idea of the proof}

In the $n$th unit cell we place a very short core of length $M_n^{-1}$.
Across that core the two weights exchange the values $M_n$ and $1$ in a
reflection-symmetric way.  A block function is constant on the core and
tapers to zero through two fixed-width collars.  Both weights are multiplied
by the tiny common factor $M_n^{-4}$ on the collars.  This common factor makes
the taper cheap but cancels from $r=w_0/w_1$.

The short core has two consequences.  First, the endpoint norm of one block
stays comparable to one.  Second, averaging over the core is uniformly
bounded on both endpoint Sobolev spaces: a function cannot change its core
average cheaply because the core is short.  These averages define a single
bounded projection on both endpoints.  Reflection makes the projected ranges
identical Hilbert sequence spaces.

For the first strict inclusion, coefficients $1/n$ put a block sum in both
endpoints, hence in the interpolation space.  The rapid exchange of the
weights forces $|(\log r)'|\asymp M_n\log M_n$ on the core, so the auxiliary
norm diverges.  For the stronger strict inclusion, coefficients $n^{-1/2}$
are square summable for the exponentially small geometric-mean block norm,
but not for the projected endpoint norm.  The common projection then rules
out membership even in $X_0+X_1$.

\section{Construction}

For $n\geq1$ set
\[
 C_n=(n,n+1),\quad z_n=n+\tfrac12,\quad
 M_n=16^n,\quad L_n=M_n^{-1},\quad h=\tfrac18,
\]
and let
\[
 a_n=z_n-\tfrac{L_n}{2},\qquad b_n=z_n+\tfrac{L_n}{2},
 \qquad I_n=(a_n,b_n).
\]
Define a continuous piecewise affine profile $s_n:C_n\to[0,1]$ by
\begin{equation}\label{eq:profile}
s_n(x)=
\begin{cases}
\frac12,&x\notin(a_n-h,b_n+h),\\
\frac12+\dfrac{x-(a_n-h)}{2h},&a_n-h<x<a_n,\\
1,&a_n\leq x\leq a_n+L_n/3,\\
2-3\dfrac{x-a_n}{L_n},&a_n+L_n/3<x<a_n+2L_n/3,\\
0,&a_n+2L_n/3\leq x\leq b_n,\\
\dfrac{x-b_n}{2h},&b_n<x<b_n+h.
\end{cases}
\end{equation}
Let the common factor be
\[
 q_n(x)=
 \begin{cases}1,&x\in I_n,\\ M_n^{-4},&x\in C_n\setminus I_n,
 \end{cases}
\]
and, for $x\in C_n$, put
\begin{equation}\label{eq:weights}
 w_0(x)=q_n(x)M_n^{s_n(x)},\qquad
 w_1(x)=q_n(x)M_n^{1-s_n(x)}.
\end{equation}
At points outside the union of the open cells, set $w_0=w_1=1$.

The $n$th block is the symmetric tent
\begin{equation}\label{eq:block}
\beta_n(x)=
\begin{cases}
0,&x\notin(a_n-h,b_n+h),\\
\dfrac{x-(a_n-h)}h,&a_n-h<x<a_n,\\
1,&a_n\leq x\leq b_n,\\
\dfrac{b_n+h-x}h,&b_n<x<b_n+h.
\end{cases}
\end{equation}
The supports of the $\beta_n$ are pairwise disjoint.

\begin{lemma}[admissibility]\label{lem:admissible}
The functions $w_0,w_1$ are strictly positive measurable weights satisfying
the compact boundedness condition.  Their quotient $r=w_0/w_1$ is locally
Lipschitz, and on $C_n$
\[
 r=M_n^{2s_n-1},\qquad w_{1/2}=\sqrt{w_0w_1}=q_nM_n^{1/2}.
\]
\end{lemma}

\begin{proof}
Every compact subset of $\R$ meets only finitely many cells.  On each of those
cells each weight assumes only finite positive values bounded above and away
from zero, so compact boundedness follows.  The common factor cancels from
the quotient.  Since $s_n=1/2$ near both endpoints of $C_n$, the cellwise
quotient joins the constant value $1$ across cell boundaries.  It is
piecewise Lipschitz in each cell; on a compact set only finitely many pieces
occur.  This proves local Lipschitz continuity.  The formula for $w_{1/2}$ is
immediate from \eqref{eq:weights}.
\end{proof}

Let $R_nx=2z_n-x$.  Directly from the construction,
\begin{equation}\label{eq:reflection}
 s_n(R_nx)=1-s_n(x),\qquad q_n(R_nx)=q_n(x),\qquad
 \beta_n(R_nx)=\beta_n(x).
\end{equation}

\section{The common complemented block subspace}

Define
\[
 B_{j,n}=\|\beta_n\|_{X_j}^2
 =\int_{C_n}w_j\bigl(|\beta_n|^2+|\beta_n'|^2\bigr),
 \qquad j=0,1.
\]

\begin{lemma}[block estimates]\label{lem:block-estimates}
For every $n$,
\begin{equation}\label{eq:B-bounds}
 B_{0,n}=B_{1,n}=:B_n,qquad
 \frac13\leq B_n\leq1+\frac{65}{4}M_n^{-3}<1.005.
\end{equation}
Moreover, if
\[
 D_n=\|\beta_n\|_{W^{1,2}(\R,w_{1/2})}^2,
\]
then
\begin{equation}\label{eq:D-bound}
 D_n\leq M_n^{-1/2}+\frac{65}{4}M_n^{-7/2}.
\end{equation}
\end{lemma}

\begin{proof}
Equation~\eqref{eq:reflection} exchanges $w_0$ and $w_1$ and fixes
$|\beta_n|^2+|\beta_n'|^2$, so $B_{0,n}=B_{1,n}$.  On the core,
$\beta_n=1$, $\beta_n'=0$, and an exact calculation gives
\begin{align*}
 \int_{I_n}w_0
 &=L_n\left(\frac{M_n}{3}
      +\frac{M_n-1}{3\log M_n}+\frac13\right)\\
 &=\frac13+\frac{M_n-1}{3M_n\log M_n}+\frac1{3M_n}.
\end{align*}
This lies between $1/3$ and $1$.  The same holds for $w_1$ by reflection.
On the two collars, $w_j\leq M_n^{-3}$, $|\beta_n|\leq1$, and
$|\beta_n'|=h^{-1}=8$.  Their total contribution is at most
\[
 2h(1+h^{-2})M_n^{-3}=\frac{65}{4}M_n^{-3},
\]
which proves \eqref{eq:B-bounds}.  For $w_{1/2}$, the core contribution is
$L_nM_n^{1/2}=M_n^{-1/2}$ and the collar weight is $M_n^{-7/2}$; the same
calculation gives \eqref{eq:D-bound}.
\end{proof}

For a locally Sobolev function $f$, let
\[
 m_n(f)=\frac1{L_n}\int_{I_n}f(x)\,dx,
 \qquad Pf=\sum_{n\geq1}m_n(f)\beta_n.
\]

\begin{lemma}[common projection]\label{lem:projection}
The same formula $P$ defines a bounded projection on both $X_0$ and $X_1$.
Its range in either endpoint is
\begin{equation}\label{eq:range}
 \mathcal R=\left\{\sum_{n\geq1}c_n\beta_n:
             \sum_{n\geq1}B_n|c_n|^2<\infty\right\}.
\end{equation}
In particular, the two endpoint range spaces coincide isometrically.
\end{lemma}

\begin{proof}
Fix $j\in\{0,1\}$.  Let $H_{j,n}$ be the third of $I_n$ on which $w_j=M_n$:
the left third for $j=0$ and the right third for $j=1$.  Weighted Sobolev
functions are absolutely continuous on the compact core.  Comparing the two
averages and using the fundamental theorem of calculus gives
\[
 |m_n(f)|\leq\left|\frac3{L_n}\int_{H_{j,n}}f\right|
              +\int_{I_n}|f'|.
\]
Cauchy--Schwarz, $M_nL_n=1$, and $w_j\geq1$ on $I_n$ yield
\begin{align*}
 |m_n(f)|^2
 &\leq \frac6{M_nL_n}\int_{H_{j,n}}w_j|f|^2
       +2L_n\int_{I_n}w_j|f'|^2\\
 &\leq6\int_{I_n}w_j(|f|^2+|f'|^2).
\end{align*}
The supports of the blocks are disjoint, so Lemma~\ref{lem:block-estimates}
implies
\[
 \|Pf\|_{X_j}^2
 =\sum_nB_n|m_n(f)|^2\leq 6\sup_nB_n\,\|f\|_{X_j}^2.
\]
Since $m_n(\beta_k)=\delta_{nk}$, $P^2=P$.  Disjointness and
$B_{0,n}=B_{1,n}=B_n$ now give exactly \eqref{eq:range} at both endpoints.
\end{proof}

\section{Both strict inclusions}

\begin{theorem}\label{thm:main}
For the weights \eqref{eq:weights},
\[
 \Waux^2(\R,1/2,r)
 \subsetneq [X_0,X_1]_{1/2}
 \subsetneq W^{1,2}(\R,w_{1/2}).
\]
Consequently, the answer to the first question in Section~8.3 of
\cite{CE} is affirmative, and the unrestricted identity
\[
 [W^{1,2}(\R,w_0),W^{1,2}(\R,w_1)]_{1/2}
 =W^{1,2}(\R,\sqrt{w_0w_1})
\]
is false even when both weights are compact-bounded and their quotient is
locally Lipschitz.
\end{theorem}

\begin{proof}
For the first strictness, set
\[
 g=\sum_{n\geq1}\frac1n\beta_n.
\]
By \eqref{eq:B-bounds},
$\|g\|_{X_j}^2=\sum_nB_n/n^2<\infty$ for $j=0,1$.  Hence
$g\in X_0\cap X_1\subset[X_0,X_1]_{1/2}$.

On the middle third of $I_n$, $\beta_n=1$ and
\[
 |(\log r)'|=\frac{6\log M_n}{L_n},\qquad
 w_{1/2}=M_n^{1/2}.
\]
Therefore that third alone contributes
\begin{equation}\label{eq:aux-divergence}
 \frac1{n^2}\frac{36(\log M_n)^2}{L_n^2},
 M_n^{1/2}\frac{L_n}{3}
 =\frac{12}{n^2}(\log M_n)^2M_n^{3/2}
\end{equation}
to
$\int |g|^2|(\log r)'|^2w_{1/2}$.  These terms do not even tend to zero, so
$g\notin\Waux^2(\R,1/2,r)$.  Lemma~\ref{lem:admissible} permits us to apply
the source inclusion \eqref{eq:source-inclusion}; hence the first inclusion is
strict.

For the second strictness, define
\[
 f=\sum_{n\geq1}\frac1{\sqrt n}\beta_n.
\]
By \eqref{eq:D-bound},
\[
 \|f\|_{W^{1,2}(w_{1/2})}^2
 \leq\sum_{n\geq1}\frac1n
 \left(M_n^{-1/2}+\frac{65}{4}M_n^{-7/2}\right)<\infty.
\]
Thus $f\in W^{1,2}(\R,w_{1/2})$.

Suppose, toward a contradiction, that $f\in X_0+X_1$.  The common formula
for $P$, bounded on each endpoint, defines a bounded map from the algebraic
sum into $\mathcal R+\mathcal R=\mathcal R$.  But $Pf=f$, because every block
is identically one on its own averaging core.  Hence \eqref{eq:range} would
give
\[
 \sum_{n\geq1}\frac{B_n}{n}<\infty,
\]
contradicting $B_n\geq1/3$.  Therefore $f\notin X_0+X_1$.  Every complex
interpolation space is an intermediate space contained in the endpoint sum,
so $f\notin[X_0,X_1]_{1/2}$.  The second inclusion in
\eqref{eq:source-inclusion} is strict.
\end{proof}

\section{Verification, limitations, and review focus}

The companion script \texttt{code/verify\_blocks.py} checks the reflection
identities on 12,006 sample points, the exact core-mass formula, the uniform
endpoint and projection bounds, the geometric-mean decay, and finite partial
sums for both witnesses.  It is a sanity check; the proof above uses the
displayed analytic inequalities.

A bounded novelty search through 13 August 2026 used the run's four cheap
indexes, local arXiv full-text sources and citation traces, and exact web/arXiv
queries for the source question and close variants.  Later related work
included arXiv:2204.11583 on density and power weights on the half-line and
arXiv:2503.14636 on complex interpolation of power-weighted Sobolev spaces
with boundary conditions.  Neither contained this block construction, an
answer to the strict-inclusion question, or a general counterexample of this
form.  This is evidence only, not a claim of exhaustive novelty.

The construction deliberately uses measurable weights with jumps in the
common factor $q_n$; this is allowed by the source's compact boundedness
hypothesis.  The quotient $w_0/w_1$, which is the quantity required to be
locally Lipschitz, is continuous and locally Lipschitz.  Human review should
focus on the boundedness and common-range assertions for $P$, and on the
application of Theorem~1.21 of \cite{CE} to pass from the full auxiliary space
to the interpolation space on $\R$.

\begin{thebibliography}{9}
\bibitem{CE}
M.~Cwikel and A.~Einav,
\emph{Interpolation of weighted Sobolev spaces},
J. Funct. Anal. 277 (2019), no.~7, 2381--2441;
\href{https://arxiv.org/abs/1807.10790}{arXiv:1807.10790}.
\end{thebibliography}

\end{document}
